\section{Auxiliary supersingular norm and order calculations}
\label{sec:auxiliary-supersingular}

The following independent calculations were developed while testing the small-prime obstruction. They are not needed for the final all-prime minimum V06, whose proof uses the stronger odd-component cancellation argument. We include their full proofs because the norm estimate and the actual order embedding also clarify which finite-characteristic transport statements are available without lifting effective cycles.

\begin{proposition}[C30: the commutator norm gap]
Let $E/\overline{\F}_p$ be supersingular with a specified embedding $\mathcal O_K=\Z[\omega]\subset\End(E)$, where $K=\Q(\omega)$ and $\omega^2+\omega+1=0$. Let $n$ be the reduced norm on $D=\End^0(E)$, and let $n_\perp(u)$ be the norm of the component perpendicular to $K$. Then
\[
n([\omega,u])=3n_\perp(u).
\]
For every actual endomorphism $u\notin K$, this is a positive integer divisible by $p$. Consequently
\[
n(u)\geq n_\perp(u)\geq p/3,
\qquad n(u)\geq\lceil p/3\rceil.
\]
In particular every endomorphism of norm at most three lies in $\mathcal O_K$ if $p\geq11$.
\end{proposition}

\begin{proof}
The supersingular endomorphism algebra is the quaternion division algebra ramified at $p$ and infinity, and the reduced norm of an endomorphism is its degree \cite[II, Proposition 6.4, p.~71; Theorem 6.9, p.~73]{MilneEC}. Write
\[
D=K\oplus Kj,\qquad jk=\bar k j,\qquad j^2=-c,
\]
where $c>0$, and set $s=\omega-\bar\omega$, so $s^2=-3$. For $u=\alpha+\beta j$ one has
\[
n(u)=\Nm(\alpha)+c\Nm(\beta),\quad
n_\perp(u)=c\Nm(\beta),\quad
[\omega,u]=s\beta j.
\]
Taking norms proves the displayed factor three.

If $u$ is an actual endomorphism, the commutator is an actual endomorphism. Its differential at the identity vanishes because the tangent space of $E$ has dimension one, so the two induced scalar maps commute. Translation then shows that its differential vanishes everywhere. If $u\notin K$, the formula above shows that the commutator is nonzero. It is therefore an inseparable isogeny of elliptic curves, whose degree is divisible by $p$; see the function-field description in \cite[V, \S7, p.~216]{MilneEC}. This proves the divisibility and the inequalities. Finally an element of $K\cap\End(E)$ is an algebraic integer, and $\mathcal O_K$ is the full ring of integers of $K$.
\end{proof}

\begin{proposition}[C31: the low-degree graph entry budget]
Let $h:E^2\to E^2$ be an isogeny with $d=\deg h\in\{1,2\}$, and suppose its graph has product-polarization degree at most fourteen. Then every entry of $h$ has norm at most three. In characteristic $p\geq11$, all its entries belong to $\mathcal O_K$.
\end{proposition}

\begin{proof}
Write $\lambda=\theta_1+\theta_2$ on $E^2$ and $T=\sum_{i,j}n(h_{ij})$. On the graph the ambient polarization pulls back to $\lambda+h^*\lambda$. Since
\[
\lambda^2=2,\qquad (h^*\lambda)^2=2d,
\qquad \lambda\,h^*\lambda=T,
\]
its degree is $2(1+d+T)$. The middle identity follows from the projection formula; the last follows by restricting the two coordinate pullback divisors to the coordinate elliptic curves, with their inseparable multiplicities included. Hence $T\leq5$ for $d=1$ and $T\leq4$ for $d=2$.

If all four entries are nonzero, their norms are positive integers, so each is at most two in the first case and at most one in the second. If an entry is zero, row and column permutations make the matrix triangular. Its degree is the product of the norms of its two nonzero diagonal entries. For $d=1$, both norms are one, leaving at most three for the off-diagonal entry. For $d=2$, their norms are one and two, leaving at most one. This covers every zero pattern of an isogeny matrix. Proposition C30 gives the final assertion.
\end{proof}

\begin{proposition}[C32: an actual characteristic-five order]
Let
\[
E/\F_5:\ y^2=x^3+1,\qquad
\omega(x,y)=(\zeta x,y),\quad \zeta^2+\zeta+1=0.
\]
Let $F$ be the $5$-power Frobenius, put $\pi=-F$ and $s=\omega-\bar\omega$. There is an actual endomorphism
\[
q=(1+\pi)s^{-1},\qquad qs=1+\pi,
\]
where division is on the right. It satisfies
\[
q^2=-2,\qquad q\omega=\bar\omega q+1,
\qquad \End(E)=\mathcal O_K\oplus\mathcal O_Kq.
\]
For $\alpha,\beta\in\mathcal O_K$,
\begin{align*}
\operatorname{pr}_K(\alpha+\beta q)&=\alpha+\beta/s,\\
n_\perp(\alpha+\beta q)&=\frac53\Nm(\beta),\\
n(\alpha+\beta q)&=\Nm(\alpha+\beta/s)+\frac53\Nm(\beta).
\end{align*}
\end{proposition}

\begin{proof}
The five affine $\F_5$-points are $(0,1),(0,4),(2,2),(2,3),(4,0)$. With the point at infinity this gives $\#E(\F_5)=6$, hence Frobenius trace zero and $\pi^2=-5$. Also $\zeta^5=\zeta^2$, so $\pi\omega=\bar\omega\pi$. The algebra is noncommutative and thus is the supersingular quaternion algebra of discriminant five by \cite[II, Theorem 6.9]{MilneEC}.

Since $s=\omega(1-\omega)$, its kernel is
\[
E[s]=\{O,(0,1),(0,-1)\}.
\]
It is etale of order three, and $F$ fixes it pointwise. Therefore $1+\pi=1-F$ kills this group scheme. Descent along the quotient isogeny $s:E\to E/E[s]\simeq E$ gives an actual endomorphism $q$ with $qs=1+\pi$. This proves the inclusion in the geometric endomorphism ring before any maximal-order argument is used.

The relation $\pi s=-s\pi$ gives $q=s^{-1}-s^{-1}\pi$. Direct multiplication now yields
\[
q^2=-2,\qquad q\omega-\bar\omega q=s^{-1}(\omega-\bar\omega)=1.
\]
Its reduced trace is zero and its norm is two. The lattice $\mathcal O=\mathcal O_K\oplus\mathcal O_Kq$ is closed under multiplication by these relations and lies inside $\End(E)$. In the ordered basis $(1,\omega,q,\omega q)$, the reduced-trace pairing $\operatorname{trd}(xy)$ has matrix
\[
\begin{pmatrix}
2&-1&0&1\\
-1&-1&1&-1\\
0&1&-4&2\\
1&-1&2&-3
\end{pmatrix},
\qquad \det=-25.
\]
Thus $\mathcal O$ has reduced discriminant five. By \cite[Theorem 15.5.5, p.~241]{Voight} it is maximal, since its reduced discriminant equals that of the quaternion algebra. The already proved inclusion into the order $\End(E)$ must therefore be equality. No class-number computation or uniqueness assertion about conjugacy classes of maximal orders is required.

Finally the $K$ component of $q=s^{-1}-s^{-1}\pi$ is $s^{-1}$, while the perpendicular component has norm $5/3$. Orthogonality gives the three asserted formulas. All geometric $j=0$ curves in characteristic five are isomorphic to this model; transporting the marking replaces $\omega$ by $\omega$ or $\bar\omega$, so the formula holds after this explicit choice of marking. It does not assert descent of the entire endomorphism ring to $\F_5$.
\end{proof}

The sign and order of division in C32 are essential. If $\pi=+F$ is used on this model, $1+\pi$ acts as multiplication by two on $E[s]$, so the same right quotient is not integral. Left division is a different expression and cannot be substituted silently.

\paragraph{C33: verification scope.}
The identifier C33 in the research memo records an independent complete proof check of V01--V04, subsequently extended to V05--V06. It is a check event, not an additional mathematical theorem. The full primitive proofs appear in their own section. In particular the final all-prime result does not depend on C30--C32 or an unproved transfer of complex positivity to arbitrary effective cycles in characteristic $p$.
